% ===================================================== BÖLÜM 10
\section{Hedef Kimlik Kartları}
\label{sec:hedefler}

\subsection*{WASP-12\,b — kalibrasyon hedefi}
\begin{tabularx}{\textwidth}{@{}l>{\raggedright\arraybackslash}X}
Konum (J2000) & RA 97,6366$^\circ$, Dec $+29{,}672^\circ$; $V = 11{,}6$; F türü yıldız, 6300 K\\
Gezegen & $M_p = 1{,}47$ M$_\mathrm{J}$, $R_p = 1{,}9$ R$_\mathrm{J}$, $P = 1{,}0914$ gün, $a = 0{,}0234$ AB\\
Literatür & Maciejewski vd. 2016 (ilk tespit); Patra vd. 2017; Yee vd. 2020 ($-29\pm2$, Q$'_\star=1{,}8\times10^5$);
            Wang vd. 2024 ($-30{,}2\pm0{,}9$); Winn \& Stefánsson 2025 (LTTE dışlandı)\\
Veri & I\&W 2022 (107 lit.), Wang TESS (82), yeni TESS s71–72 (38), Yee 2020 (139 geçiş + 19 örtülme, 72 RV)\\
Ön sonuç & $\dot P = -31{,}05\pm0{,}82$\msyil; Rømer dışlandı; apsis dışlandı; gelgit en olası\\
Başarı ölçütü & Literatürle 2$\sigma$ içinde uyum $\Rightarrow$ yöntem doğrulandı
\end{tabularx}

\subsection*{WASP-4\,b — gelgit mi, ivme mi?}
\begin{tabularx}{\textwidth}{@{}l>{\raggedright\arraybackslash}X}
Konum & RA 353,563$^\circ$, Dec $-42{,}06^\circ$; $V = 12{,}5$; G7V, 5400 K\\
Gezegen & $M_p = 1{,}2$ M$_\mathrm{J}$, $P = 1{,}3382$ gün\\
Literatür & Bouma vd. 2019, 2020 (RV ivmesi, $-5{,}94$ öngörü); Baştürk vd. 2025 (“bozuluyor”, Q$'\sim8\times10^4$);
            Winn \& Stefánsson 2025 (“bozulmuyor”, $\sim$8 AB'de $\sim$7 M$_\mathrm{J}$ eşlikçi)\\
Veri & I\&W (75 lit.), Wang TESS (47), yeni TESS s36, 69, 96, 103, 105, 106 (75), Bouma 48 RV\\
Ön sonuç & $\dot P = -5{,}11\pm0{,}30$; Rømer öngörüsü $-5{,}95\pm0{,}32$ (1,9$\sigma$); kübik $\Delta$BIC = 13\\
Proje sorusu & $|\dot P|$'nin zamanla küçülmesi LTTE'yi mi gösteriyor? Kalan küçük fark gelgit için yer bırakıyor mu?
\end{tabularx}

\subsection*{HAT-P-7\,b — dönemi uzayan sistem}
\begin{tabularx}{\textwidth}{@{}l>{\raggedright\arraybackslash}X}
Konum & RA 292,247$^\circ$, Dec $+47{,}97^\circ$; $V = 10{,}5$; F6, 6400 K; Kepler alanında\\
Gezegen & $M_p = 1{,}84$ M$_\mathrm{J}$, $P = 2{,}2047$ gün; kutupsal/ters yörünge\\
Literatür & Winn vd. 2009 (RV eğilimi $21{,}5\pm2{,}6$ m/s/yıl); Narita vd. 2012 (HAT-P-7\,c; görüntülemede yıldız eşlikçi);
            Knutson vd. 2014 (eğilim tepe yaptı); Wang vd. 2024 ($+18{,}3\pm4{,}1$)\\
Veri & I\&W (11 lit., 9'u zaman sistemi belirsiz), Wang TESS (66), yeni TESS s74–86 (40), Knutson 43 RV\\
Ön sonuç & $\dot P = +28{,}1\pm3{,}3$; Rømer öngörüsü $+13{,}1\pm1{,}2$ (kısmen); gelgit dışlandı\\
Proje sorusu & İvmenin değişimi (eğrilik) ve zaman sistemi düzeltmesi farkı kapatıyor mu?
\end{tabularx}

\subsection*{TrES-1\,b — ne gelgit ne Rømer}
\begin{tabularx}{\textwidth}{@{}l>{\raggedright\arraybackslash}X}
Konum & RA 286,041$^\circ$, Dec $+36{,}63^\circ$; $V = 11{,}4$; K0V, 5200 K\\
Gezegen & $M_p = 0{,}84$ M$_\mathrm{J}$, $P = 3{,}0301$ gün\\
Literatür & Alonso vd. 2004 (keşif); Adams vd. 2024 (Q$'$ çok küçük $\Rightarrow$ bozulma değil); Hagey vd. 2025 (apsis,
            olası TrES-1\,c); Kutluay vd. 2026 ($-14{,}9$)\\
Veri & I\&W (41 lit.), Wang TESS (34), yeni TESS s74–86 (21), HARPS-N 15 RV (Bonomo vd. 2017), DACE 33 RV\\
Ön sonuç & $\dot P = -12{,}7\pm1{,}9$; gereken Q$'_\star = 8\times10^2$ (gelgit olası değil); Rømer öngörüsü $-2{,}5\pm2{,}2$ (3,5$\sigma$ fark: dışlandı); Roche 3,1\\
Dikkat & Erken (2004–2006) noktalar çok saçılmış; kaynak çıkarma testi şart
\end{tabularx}

% ===================================================== BÖLÜM 11
\section{Başvuruya Kadar Gün Gün Plan (6–11 Kasım 2026)}
\label{sec:basvuru}

\begin{dikkat}[Tarihler]
Son başvuru: \textbf{11 Kasım 2026 Çarşamba, 17:30} (TÜBİTAK BİDEB duyurusu; sistem 17:30'da kapanır).
Kuruluş yetkilisi e-imzası: 16–23 Kasım 2026. Danışman onayının da sistemde zamanında verilmesi gerekir:
\textbf{başvuruyu en geç 10 Kasım akşamı gönder}, danışmana onay için bir gün bırak.
\end{dikkat}

\begin{center}
\small
\begin{tabularx}{\textwidth}{l >{\raggedright\arraybackslash}X >{\raggedright\arraybackslash}X}
\toprule
Gün & Yapılacak & Çıktı\\
\midrule
6 Kas (bugün) & Kılavuzun Bölüm~\ref{sec:secim}'ini oku; hocana seçimi (4 hedef) ve gerekçe paragrafını göster, onay al. & Hoca onayı\\
7 Kas & Kısım I'deki metinleri Word formuna yapıştır (Arial 9). Danışman adını yaz. & Dolu form taslağı\\
8 Kas & Bölüm 5'e Şekil~1, Şekil~2, Tablo~1, Tablo~2'yi ekle (\texttt{sekiller/}). Sayfa sayısını kontrol et (ekler hariç $\le$20). & Tam form\\
9 Kas & Formu hocana gönder; düzeltmeleri işle. & Hoca onaylı form\\
10 Kas & TÜBİTAK sistemine (TYBS) yükle ve gönder. Danışmana “onay verebilirsiniz” mesajı at. & Gönderilmiş başvuru\\
11 Kas & Yedek gün (sistem sorunu, onay gecikmesi). & --\\
Başvurudan sonra & Kodları Colab'da bir kez çalıştır (Bölüm~\ref{sec:kod}); sayıları bu belgedekilerle karşılaştır. Mülakat/soru için hazırlık. & Kodu anlatabiliyor olmak\\
\bottomrule
\end{tabularx}
\end{center}

\subsection{Forma konacak şekil ve tablolar (Bölüm 5 “Belirtmek istediğiniz diğer konular”)}
\begin{enumerate}[leftmargin=*]
\item \textbf{Şekil 1}: WASP-12\,b O$-$C diyagramı (yeni TESS zamanları kırmızı). Altyazı: “Ön çalışma: Ivshina \& Winn
      (2022) ve Wang vd. (2024) zamanlarına bu çalışmada TESS sektör 71–72'den ölçülen 38 yeni geçiş zamanı eklenmiştir.
      Kuadratik model $\dot P = -31{,}1\pm0{,}8$ ms/yıl verir; literatürdeki değerle ($-29\pm2$; Yee vd. 2020) uyumludur.”
\item \textbf{Şekil 2}: Dört hedefin O$-$C diyagramları (Şekil~\ref{fig:dort}).
\item \textbf{Tablo 1}: 14 adaydan seçim (Kısım I'de hazır).
\item \textbf{Tablo 2}: Rømer testi ve gereken Q$'_\star$ (Kısım I'de hazır).
\end{enumerate}
Word'de şekil eklemek: Ekle $\rightarrow$ Resimler $\rightarrow$ \texttt{sekiller/sekil1\_WASP-12b\_OC.png};
genişliği sayfa genişliğine ayarla; altına “Şekil 1.~\dots” yaz.

% ===================================================== BÖLÜM 12
\section{Proje Süresince 12 Aylık Plan}
\label{sec:plan}
Ay 1, TÜBİTAK'ın bildireceği başlangıç ayıdır. 2209-A kuralı gereği literatür taraması, rapor yazımı ve sonuç paylaşımı
ayrı iş paketi olarak gösterilmez; iş paketlerinin içinde yürür.

\begin{center}
\small
\begin{tabularx}{\textwidth}{l l >{\raggedright\arraybackslash}X >{\raggedright\arraybackslash}X r}
\toprule
İP & Ay & İş & Başarı ölçütü & Katkı (\%)\\
\midrule
1 & 1–3 & Katalog ve literatür geçiş zamanlarının derlenmesi (I\&W, Wang, ETD/VarAstro, ExoClock, hedefe özel makaleler); BJD$_\mathrm{TDB}$'ye dönüşüm; kaynak etiketleme &
  4 hedef için her birinde $\ge$70 zaman; zaman sistemi belirsiz kaynak oranı $\le$\%10 ve ayrı işaretli & 20\\
2 & 2–5 & Wang sonrası TESS sektörlerinden orta geçiş zamanlarının ölçülmesi & $\ge$150 yeni zaman; ortanca hata $\le$60 s; Wang ile ortak geçişlerde ortanca fark $\le$1$\sigma$ & 20\\
3 & 4–7 & O$-$C diyagramları; doğrusal/kuadratik/kübik model; BIC ile model seçimi &
  4 hedef için $\dot P\pm\sigma$, $\Delta$BIC; WASP-12'de literatürle 2$\sigma$ uyum & 25\\
4 & 6–11 & Fiziksel köken: arşiv RV ile Rømer testi; Q$'_\star$; apsis (örtülme, $e$); Applegate (etkinlik, periyodogram); dinamik 3. cisim; Roche &
  4 hedef $\times$ 6 etki = 24 hücrelik karar tablosunun her hücresinde gerekçeli hüküm; her hedef için sayısal RV öngörüsü & 35\\
\bottomrule
\end{tabularx}
\end{center}

% ===================================================== BÖLÜM 13
\section{Değerlendirme: Bu Kurgu 2209-A İçin Yeterli mi?}
\label{sec:degerlendirme}

\subsection{Kısa cevap: Evet, yeterli.}
2209-A bir \emph{lisans öğrencisi} araştırma programıdır. Hakemler devrim değil, \textbf{net bir araştırma sorusu,
uygulanabilir bir yöntem, ölçülebilir başarı ölçütleri ve gerçekçi bir takvim} arar. Bu kurgunun güçlü yanları:
\begin{itemize}[leftmargin=*]
\item \textbf{Net soru}: “Dönem değişiyor mu? Değişiyorsa neden?”
\item \textbf{Gerekçeli örneklem}: Sayıyı ölçütlerin belirlediği, 14'ten 4'e inen bir huni (Bölüm~\ref{sec:secim}).
\item \textbf{Yöntem doğrulaması içeride}: WASP-12\,b bilinen sonucu veriyor (pozitif kontrol).
\item \textbf{Ön çalışma gerçek}: Şekiller ve sayılar var; “yapacağız” değil, “başladık, işe yarıyor”.
\item \textbf{Özgün katkı}: (1) Wang vd. sonrası 174 yeni TESS geçiş zamanı (2–5 yıl daha uzun zaman tabanı);
      (2) her hedef için altı etkinin \emph{sistematik} değerlendirildiği karar tablosu; (3) arşiv RV'siyle Rømer testi;
      (4) her hedef için test edilebilir RV öngörüsü. WASP-4'teki “$|\dot P|$ küçülüyor” bulgusu gibi gerçek sonuçlar
      çıkabiliyor.
\item \textbf{Riskler düşük}: Bütçe 0 TL, gözlem zamanına ve hava koşuluna bağımlılık yok, tüm veri açık.
\end{itemize}

\subsection{Değiştirmeni önerdiğim şeyler}
\begin{enumerate}[leftmargin=*]
\item \textbf{Eski raporun kapsamını formdan sil}: Gaia/HGCA, 26 aday, Q$'_\star$ Monte Carlo, ATA50 gözlemi. Bunlar
      hocanın istemediği ve 12 ayda zorlanacağın işler.
\item \textbf{“Yeni RV gözlemi yapılmayacak, yalnızca yayımlanmış arşiv RV'leri kullanılacak” cümlesini formda açıkça yaz}.
      Hakem “tayfçekeriniz var mı?” diye düşünmesin.
\item \textbf{Yeni TESS sektörlerini kesinlikle koru}: Projeyi “başkasının tablosunu yeniden çizmek”ten ayıran şey budur.
      Hocanın “TESS verisini dahil et” sözünün tam karşılığıdır.
\item \textbf{Hipotezleri “sınamak” diliyle yaz}, “kanıtlamak” değil. Hipotez reddedilse de sonuç yayımlanabilir.
\item \textbf{“İlk kez” deme}: WASP-12 ve WASP-4 çok çalışılmış. Bunları “kalibrasyon” ve “tartışmalı vaka” olarak sun.
      Bu dürüstlük hakemin gözünde artıdır.
\item \textbf{Başarı ölçütleri sayı olsun} (Bölüm~\ref{sec:plan}); toplamları 100.
\end{enumerate}

% ===================================================== BÖLÜM 14
\section{Hocanın ya da Jürinin Sorabileceği Sorular}
\label{sec:sorular}
\begin{description}[leftmargin=0pt,style=nextline]
\item[“LOOCV'yi geçmek ne demek?”] Veriden hangi tek ölçümü çıkarırsam çıkarayım dönem değişimi hâlâ 3$\sigma$ ve
  $\Delta$BIC $>10$ düzeyinde anlamlı kalıyor demek. Sonucun tek bir hatalı ölçüme bağlı olmadığını gösterir; ama fiziksel
  sebebi söylemez ve aynı makaleden gelen bir grup hatalı noktaya karşı korumaz.
\item[“Neden 4 hedef? Neden bunlar?”] Sayıyı biz seçmedik. Sonuçlara bakmadan yazdığımız ölçütleri uyguladık: Wang vd.'nin
  LOOCV listesindeki 14 aday içinden Wang vd.'ye göre en az 30 TESS geçişi, Rømer öngörüsünü gözlenen sinyalin üçte birinden hassas verecek arşiv RV'si.
  Bunları sağlayan dört hedef kaldı (Bölüm~\ref{sec:secim}'teki hazır paragraf).
\item[“Neden en güçlü $\Delta$BIC'li adayları seçmedin?”] $\Delta$BIC veri sayısıyla büyür ve iyimser hata çubuklarıyla
  şişer; tek bir aykırı nokta onu ciddi biçimde değiştirebilir (Wang vd. 2024). Bu yüzden güvenilirlik ölçütü olarak
  kullanmadık; anlamlılık için $|\dot P|/\sigma$, model seçimi için $\Delta$BIC $>10$ kullandık.
\item[“Gözlem yapmadan bu nasıl özgün olur?”] Wang vd.'nin kullanmadığı 2022–2025 TESS sektörlerinden 174 yeni geçiş
  zamanını kendim ölçüyorum. Altı etkiyi aynı karar ağacıyla sistematik olarak sınıyorum.
\item[“Gelgit ile Rømer'i nasıl ayırdın?”] Zamanlama tek başına ayıramaz (kuadratik terimler matematiksel olarak aynı).
  Arşivdeki RV eğiminden $\dot P_\text{Rømer} = P\dot v_r/c$'yi hesaplayıp ölçtüğüm $\dot P$ ile karşılaştırdım.
\item[“RV yeterli olmasaydı?”] B planı: literatürdeki eğimler, Q$'_\star$ olabilirliği, eğrilik/periyodiklik testleri ve
  “Rømer olsaydı RV şu hızla değişmeliydi” öngörüsü.
\item[“Q$'_\star$ nedir?”] Yıldızın gelgit enerjisini ne kadar verimli soğurduğunun ölçüsü; küçükse bozulma hızlıdır.
  $10^5$–$10^7$ beklenir; $10^3$ gerektiren bir değişim gelgitle açıklanamaz.
\item[“Applegate'i nasıl dışlarsın?”] F türü (etkinliği zayıf) yıldızlarda olası değildir. G/K yıldızlarda periyodik
  sinyal ve eğrilik değişimi aranır; veri birkaç döngüyü kapsamıyorsa “dışlanamadı” diye raporlanır.
\item[“Rømer ile 3. cisim aynı şey değil mi?”] İkisi de üçüncü bir cisimden gelir. Rømer \emph{geometriktir}: sistem
  hareket eder, ışık yolu değişir, yörünge aynı kalır. Dinamik etkide ise üçüncü cisim gezegenin yörüngesini kütle
  çekimiyle gerçekten bozar (kısa dönemli TTV, zorlanmış presesyon).
\item[“HAT-P-7'de dönem neden uzuyor?”] Gelgit yavaş dönen yıldızda dönemi kısaltır, uzatamaz. Bilinen eşlikçinin RV
  ivmesi artışın yaklaşık yarısını açıklıyor; kalanı ivmenin değişimi ve eski zamanların zaman sistemiyle ilgili
  olabilir. Bunu proje sırasında sınayacağım.
\end{description}

\appendix
% ===================================================== EK B
\section{Sözlük}
\begin{description}[leftmargin=0pt,style=nextline,itemsep=0.2em]
\item[Apsis hareketi (apsidal motion)] Dışmerkez yörüngenin enberi noktasının zamanla dönmesi.
\item[BIC] Bayes Bilgi Ölçütü; $\chi^2 + k\ln n$. Küçük olan model tercih edilir.
\item[BJD$_\mathrm{TDB}$] Güneş Sistemi kütle merkezine taşınmış, düzgün akan zaman. Projenin tek zaman sistemi.
\item[BTJD] TESS zamanı: BJD$_\mathrm{TDB}$ $-$ 2457000.
\item[Efemeris] Bir gezegenin geçiş zamanlarını öngören formül ($T_0$, $P$ ve varsa $\dot P$).
\item[Epok ($E$)] Referans geçişten bu yana geçen tam tur sayısı.
\item[ETD / VarAstro] Çek Astronomi Derneği'nin geçiş zamanı veritabanı ve değişen yıldız portalı.
\item[LOOCV] Birini dışarıda bırak çapraz doğrulama: sonucun tek bir ölçüme bağlı olup olmadığı testi.
\item[LTTE / Rømer etkisi] Uzak eşlikçi nedeniyle ışığın yol süresinin değişmesi.
\item[O$-$C] Gözlenen eksi hesaplanan geçiş zamanı.
\item[Örtülme (ikincil tutulma)] Gezegenin yıldızın arkasına geçmesi.
\item[Q$'_\star$] Yıldızın değiştirilmiş gelgit kalite faktörü.
\item[Roche sınırı] Gezegenin yıldızın gelgit kuvvetiyle parçalanmaya başladığı uzaklık.
\item[RV (radyal hız)] Yıldızın bakış doğrultusundaki hızı; Doppler kaymasından ölçülür.
\item[SPOC] TESS'in resmî veri işleme hattı; 2 dakikalık ışık eğrilerini üretir.
\item[TTV] Geçiş zamanı değişimi (transit timing variation).
\item[$\dot P$] Yörünge dönemi değişim hızı, ms/yıl.
\end{description}

% ===================================================== KAYNAKÇA
\section*{Kaynakça (Kılavuzda Geçen Tüm Kaynaklar)}
\addcontentsline{toc}{section}{Kaynakça}
\begin{small}
\begin{itemize}[leftmargin=*,itemsep=0.1em]
\item Adams, E. R., vd. 2024, “Doomed Worlds I: No new evidence for orbital decay in a long-term survey of 43 ultra-hot Jupiters”, arXiv:2404.07339
\item A-thano, N., vd. 2026, “Investigation of Transit Timing and an Optical Transmission Spectrum of the Hot Jupiter WASP-11 b”, arXiv:2604.05570
\item Applegate, J. H. 1992, ApJ, 385, 621
\item Baştürk, Ö., vd. 2025, “The Orbit of WASP-4 b is in Decay”, arXiv:2506.15022
\item Bouma, L. G., vd. 2019, AJ, 157, 217
\item Bouma, L. G., vd. 2020, “WASP-4 is Accelerating Toward the Earth”, ApJL, 893, L29, arXiv:2004.00637; VizieR J/ApJ/893/L29
\item Hagey, S. R., vd. 2025, “TrES-1 b: A Case Study in Detecting Secular Evolution of Exoplanet Orbits”, arXiv:2508.15075
\item Ivshina, E. S., \& Winn, J. N. 2022, ApJS, 259, 62; VizieR J/ApJS/259/62
\item Jackson, B., vd. 2016, CeMDA, 126, 227
\item Kass, R. E., \& Raftery, A. E. 1995, JASA, 90, 773
\item Knutson, H. A., vd. 2014, ApJ, 785, 126; VizieR J/ApJ/785/126
\item Kutluay, vd. 2026, “Chasing the Tides: Searching for Orbital Decay Signatures in Transit Timing Data and Tidal Models for 20 Hot Jupiters”, arXiv:2602.01446
\item Narita, N., vd. 2012, PASJ, 64, L7
\item Patra, K. C., vd. 2017, AJ, 154, 4
\item Penev, K., vd. 2018, AJ, 155, 165
\item Rothmeier, M., vd. 2025, “Doomed Worlds II: Reassessing Suggestions of Orbital Decay for TrES-5 b”, arXiv:2512.13937
\item Sodickson, N., \& Grunblatt, S. 2025, “In Search of Decay”, arXiv:2508.18355
\item Valsecchi, F., vd. 2015, ApJ, 813, 101
\item Wang, W., vd. 2024, ApJS, 270, 14, doi:10.3847/1538-4365/ad0847, arXiv:2310.17225; VizieR J/ApJS/270/14
\item Watson, C. A., \& Marsh, T. R. 2010, MNRAS, 405, 2037
\item Winn, J. N., vd. 2009, ApJL, 703, L99 (arXiv:0908.1672)
\item Winn, J. N., \& Stefánsson, G. 2025, “Orbital decay candidates reconsidered: WASP-4 b is not decaying and Kepler-1658 b is not a planet”, arXiv:2510.05229
\item Yee, S. W., vd. 2020, “The Orbit of WASP-12b is Decaying”, ApJL, 888, L5, arXiv:1911.09131; VizieR J/ApJ/888/L5
\end{itemize}
\end{small}

\end{document}
